%%%%%%%%%%%%%%%%%%%%%%%%%%%%%%%%%%%%%%%%%%%%%%%%%
%%%%%%%%%%%% chap: problem description %%%%%%%%%%
%%%%%%%%%%%%%%%%%%%%%%%%%%%%%%%%%%%%%%%%%%%%%%%%%

\chapter{Problem Description}\label{chapter:problem}

\section{From Model Output to Data-Conditioned Output}

The practical problem behind this thesis appears when a large language model is placed inside a
system that retrieves and consumes external information at inference time. In that setting, the
final answer is no longer determined only by the model parameters. It also depends on what
documents are available, how they are normalized, which fragments are selected, and how they
are injected into the prompt or intermediate context.

This changes the engineering question substantially. The task is no longer only to improve the
base model or choose a better prompt. The task becomes one of controlling and auditing a
pipeline in which external content can affect:
\begin{itemize}
    \item whether a page is mentioned at all;
    \item how early and how often its terminology appears;
    \item whether its important facts survive summarization;
    \item whether unsupported claims are introduced during rewriting;
    \item whether the resulting output is suitable for publication.
\end{itemize}

In a generative-search setting, these effects are especially visible. A user may ask a query such
as ``how to set up distributed tracing,'' and the system will respond with a synthesized answer
instead of a plain list of links. If one document has a clearer structure, stronger keyword
coverage, and better contextual support, it is more likely to influence the generated answer.
Likewise, if a document is rewritten to expose terms and facts more explicitly, its measured
visibility may increase. This is useful for optimization, but it also creates room for
over-optimization and manipulation.

\section{Core Challenges}

Designing a practical platform for this environment requires addressing several challenges at the
same time.

\subsection{Visibility Is Difficult to Observe Directly}

Commercial generative engines are typically closed systems. Their internal ranking logic,
retrieval stack, prompt construction, and answer synthesis steps are not fully observable.
Because of that, an engineering platform cannot rely solely on internal model details. It needs
proxy metrics and repeatable evaluation procedures that approximate how visible a document is
within generated answers.

\subsection{Optimization Can Damage Meaning}

The simplest way to optimize for visibility is to rewrite content aggressively around keywords,
headings, and answer-like snippets. However, this can damage the original meaning of the
source page. Numeric claims can be altered, supporting nuance can disappear, and new claims
can be invented accidentally. For a trustworthy workflow, optimization must therefore be
constrained by preserved facts, forbidden claims, and reviewable diffs.

\subsection{External Context Must Be Useful but Read-Only}

Many pages should be improved with the help of surrounding organizational knowledge. For
example, a setup guide can benefit from a concepts page and a reference page. At the same
time, that surrounding context should not be silently merged into the source page without
traceability. The system therefore needs a mechanism for attaching supporting context in a
read-only manner.

\subsection{Evaluation and Publication Should Be Separated}

Even when an optimization draft appears beneficial, it should not be published automatically.
The workflow must distinguish between:
\begin{enumerate}
    \item measuring the current state;
    \item generating a recommendation;
    \item reviewing safety checks and projected lift;
    \item approving the change;
    \item exporting the final draft.
\end{enumerate}

This separation is necessary because the thesis addresses a high-leverage influence surface.
Once external data can steer an LLM answer, every optimization step should be observable and
reversible.

\section{System Requirements}

The NewGEO platform was designed around the challenges above. The requirements can be
grouped into functional and non-functional categories.

\subsection{Functional Requirements}

The system must satisfy the following functional requirements:
\begin{itemize}
    \item create and manage projects representing monitored content properties;
    \item import or crawl pages and preserve normalized content snapshots;
    \item attach context packs containing briefs, voice rules, supporting documents, and constraints;
    \item define query clusters so evaluation is organized around realistic information needs;
    \item execute benchmark runs and persist both aggregate scores and raw run artifacts;
    \item generate recommendation bundles with rewritten drafts, diffs, rationale, and projected score deltas;
    \item block approval if critical constraint checks fail;
    \item export recommendations only after explicit approval, except for a deliberate forced-export path;
    \item support synchronous execution as well as queued background processing.
\end{itemize}

\subsection{Non-Functional Requirements}

The platform must also satisfy several non-functional requirements:
\begin{itemize}
    \item \textbf{Reproducibility}: the system should run locally without mandatory external services;
    \item \textbf{Traceability}: important transitions should be stored as explicit models, not implicit side effects;
    \item \textbf{Extensibility}: storage backends and benchmark connectors should be replaceable behind stable interfaces;
    \item \textbf{Safety}: recommendation generation should surface risks before publication;
    \item \textbf{Transparency}: scores and rewrites should be explainable to a reviewer;
    \item \textbf{Product readiness}: the workflow should be visible through a usable dashboard rather than only through scripts.
\end{itemize}

These requirements explain several implementation choices in the repository. The current
system uses deterministic heuristics instead of live engine APIs, JSON and SQLite instead of a
full database stack, and a polling worker instead of a distributed queue. These decisions reduce
infrastructure complexity while keeping the architecture open for future replacement.

\section{Operational Scenario}

The intended use of the system can be summarized through the following scenario.

\begin{enumerate}
    \item A user creates a project representing a documentation site or knowledge base.
    \item The user imports one or more pages to be monitored and optimized.
    \item For a target page, the user defines a context pack containing supporting documents, required terminology, locked facts, and forbidden claims.
    \item The user defines a query cluster that approximates the demand the page should satisfy.
    \item The system runs a benchmark to estimate how visible the page is under the tracked query cluster.
    \item The system generates a recommendation draft that restructures the content conservatively and previews the projected score lift.
    \item A reviewer inspects the rationale, diff, supporting context used, and constraint checks.
    \item If the draft is acceptable, the reviewer approves it and exports the rewritten markdown.
\end{enumerate}

This scenario is already implemented in the repository through the service layer, API routes,
seeded demo, dashboard views, and tests. It therefore serves both as a conceptual use case and
as the practical basis for the remaining chapters.

\section{Scope of the Current Prototype}

The prototype focuses on the first platform layer needed to study data-driven influence on LLM
responses. It includes ingestion, normalization, context management, heuristic benchmarking,
guarded recommendation generation, approval, export, and dashboard presentation. It does not
yet include production-grade live connectors, vector databases, robust distributed job
coordination, or direct integration with real content management systems.

This limited scope is intentional. The thesis goal is to demonstrate a coherent and inspectable
architecture for studying and operationalizing external influence, not to claim a finished
production deployment. The next chapter presents how this scope is expressed in the system
architecture.
